\chapter{如何寫出暢銷小說}

這場講座一開始先把我們往錯的方向引去。我們先看到那間書房、那扇祕密門、那些書架，以及書如何在其中被寫出來的地方；緊接著，這個誘惑便立刻被撤回：機場休息室可以寫，飛機上可以寫，車裡也可以寫。真正的教訓不在那個房間，而在過程。從這裡開始，整場訪談便展開成一連串彼此相扣的問題：資訊如何被排列、評論如何變成診斷、為什麼修訂會被視為重新作曲、這種重新作曲要花多久，以及當一部長篇手稿忽然讓作者覺得自己根本看不下去時，究竟該怎麼辦。

\section{那個房間，以及迅速轉向方法}

這個開場之所以有用，正因為它先暫時誤導了我們。我們從場所開始，但講者拒絕讓場所承擔解釋性的重量。那間書房很舒適，甚至可能是最好的地方；但我們立刻被告知，寫作完全可以在別處進行。這樣一來，真正的主題才被清理出來：不是氛圍，而是方法。

這第一個轉向非常重要。較弱的講座會停留在作家式氛圍的層次，而這場講座沒有。它把房間當作入口，卻把過程當作內容。因此，我們幾乎立刻就準備好進入真正的轉折點：三個螢幕的配置，以及一部小說如何從一稿走到下一稿的機械。

\subsection{Question \& Answer}

\paragraph{Question.} 作者需要的是理想中的房間，還是能運作的方法？

\paragraph{Answer.} 講座的答案是：一個好的房間固然有幫助，但並不具有決定性。真正重要的是一套可重複的方法。房間可以支撐工作，卻無法解釋工作。

\section{三個螢幕與修訂的流動}

一旦房間被降格，訪談便變得技術性起來。我們看到的不是情調，而是一套流程。一個螢幕上放著第一稿；另一個螢幕上放著讀過第一稿的讀者所留下的評論；中間則是正在進行中的第二稿。這套配置足夠簡單，可以直接寫成一張圖：
\begin{equation}
\text{第一稿} \to \text{讀者評論} \to \text{第二稿}.
\end{equation}
這是整場講座真正的第一條數學脊骨。它說明，寫作並不只是產出之後再加上一點模糊反省；它是一種把材料經由回應導入重構的受控流動。

同一個點，我們還可以更尖銳地再寫一次：
\begin{equation}
\text{讀者回應} \to \text{診斷} \to \text{修訂優先順序}.
\end{equation}
一則評論的意義，不只是因為它被說出來，而在於：它是否揭露了手稿實際上正在對讀者做什麼。

\begin{definition}
所謂 \emph{修訂優先順序}，是指那些會被推到最前面的評論或診斷，因為它們指認出的是結構性的弱點，而不只是局部的個人偏好。
\end{definition}

由於沒有任何經過驗證的螢幕配置視覺證據留下來，我們只能用示意方式記錄這種安排，而不將其圖像化成具體場景：

\begin{figure}[h]
\centering
\begin{tikzpicture}
\node[draw, rectangle, minimum width=3.6cm, minimum height=1cm] (A) at (0,0) {第一稿};
\node[draw, rectangle, minimum width=3.6cm, minimum height=1cm] (B) at (5,0) {讀者評論};
\node[draw, rectangle, minimum width=3.6cm, minimum height=1cm] (C) at (10,0) {第二稿};
\draw[->] (A) -- (B);
\draw[->] (B) -- (C);
\end{tikzpicture}
\end{figure}

這張重建圖刻意保持極簡。講座要展示的不是設備崇拜，而是一種工作流程：稿子被推出去，回應被收回來，而下一稿則在這些已學到的東西所形成的壓力之下被寫出來。

\subsection{Question \& Answer}

\paragraph{Question.} 修訂時，作者面前究竟應該放哪些資訊？

\paragraph{Answer.} 講座的答案非常精確：目前的稿子、這份稿子所引發的回應，以及正在被寫出的下一稿。當這三者被明白地放在彼此關係之中，修訂就會變得清楚。

\section{從評論到診斷：節奏作為可測量的工藝問題}

接著，講座直接讓這套系統運轉給我們看。一位讀者說，他沒有像讀前兩本那樣快速地被新書吸進去，也沒有讀得那麼快。Follett 並不把這當成傷害，而是把它當成資料。於是，推論幾乎是立刻完成的：
\begin{equation}
\text{沒有迅速被吸進去} \Rightarrow \text{開場章節可能過於拖重}.
\end{equation}
這是一個工藝推理的典型動作。表面的評論被翻譯成了一個結構性問題。

接著，這個問題又被拿去和 Follett 明確提出的一條節奏經驗法則做比較：
\begin{equation}
\text{情節轉折間隔} \approx 3\text{--}4 \ \text{頁}.
\end{equation}
這條規則並不是被當成小說的普遍定理提出來的，而是被當成一種敘事牽引力的工作性測量。如果開頭轉得不夠頻繁，那麼讀者較慢才被吸進去，也就完全不是謎了。

\paragraph{Worked example.} 這套推理可以分步寫成：
\begin{enumerate}
\item 一位讀者說自己被吸引進去的速度太慢。
\item 作者拒絕把這句話聽成針對自己的侮辱。
\item 這句話被翻譯成一個結構性假設：開場可能太拖重。
\item 然後，再把這個假設拿去對照一條節奏上的工作準則。
\item 因此，這則評論便提升為高優先順序的修訂目標。
\end{enumerate}

接著，訪談者問出那個必然會來的問題：一位成名作家會不會乾脆把這種評論揮掉？回答又給了我們另一條緊湊關係：
\begin{equation}
\text{回饋} \neq \text{侮辱}, \qquad
\text{回饋} \Rightarrow \text{我能從中學到什麼？}
\end{equation}
這條等式極其重要。沒有它，整套三螢幕系統就會崩塌成一種虛榮戲法；有了它，回應才會成為儀器，而不再是冒犯。

\subsection{Question \& Answer}

\paragraph{Question.} 一則評論究竟如何變成一個結構診斷？

\paragraph{Answer.} 方法是把反應翻譯成工藝術語。``我沒有很快被吸進去'' 不再被看成純粹的口味判斷，而被轉換成關於節奏的陳述，接著再成為修訂的目標。

\section{第二稿作為重新作曲}

講座裡最強的一條主張，接著便出現了。Follett 說，當他進入第二稿時，他不會拿著第一稿去局部調整。他不是在通常意義上做 editing。他會把整本書再打一次。這個區分強烈到值得以正式方式寫下來：
\begin{equation}
\text{第二稿} \neq \text{加了註記的第一稿}, \qquad
\text{第二稿} = \text{從頭重寫的一本書}.
\end{equation}
這比一般關於修訂的建議更激進得多。第二稿不是第一個對象被修補過後的版本，而是在更充分的資訊之下重新建造出來的一個新對象。

這也正是為什麼中間那個螢幕如此重要。中間不是單純修補的地方，而是書被重新製造的地方。評論和先前頁面只是引導，它們不是某個正在被化妝改善的東西本身。

\subsection{Question \& Answer}

\paragraph{Question.} 為什麼要整本重寫，而不是單純編輯第一稿？

\paragraph{Answer.} 因為照這個觀點來看，第一稿還不是一個值得拿來打磨的正式對象。它只是材料、證據，以及診斷得以發生的契機。第二稿才是那本書在更高標準與更清楚理解之下被重新建起來的地方。

\begin{remark}
講座把這件事呈現為 Follett 自己的方法與信念。我們應當把它保存為這場訪談中的一項強工作標準，而不是把它誇大成普遍適用於所有小說的定律。
\end{remark}

\section{長篇小說的算術}

一旦方法被陳述出來，訪談者自然會提出下一個反對意見：這樣不是會花更久嗎？Follett 的回答並不閃躲，而是直接用算術：
\begin{equation}
\text{規劃} = 1\ \text{年}, \qquad
\text{第一稿} = 1\ \text{年}, \qquad
\text{第二稿} = 1\ \text{年}.
\end{equation}
因此，整個計畫的時間跨度便隱含地成為
\begin{equation}
T_{\text{書}} = T_{\text{規劃}} + T_{\text{第一稿}} + T_{\text{第二稿}}.
\end{equation}
這是整場講座最乾淨的貢獻之一。它把一種籠統的「小說要花很多時間」，拆成了一個具結構的時間式樣。

接下來的質疑也就順理成章：若第二稿如此昂貴，它真的有必要嗎？Follett 的答案是肯定的。出版社有時也許會希望直接出版第一稿，但在他看來，那代表的是較低的標準。他自己的結論可以寫成：
\begin{equation}
\text{從頭重寫} \Rightarrow \text{更好的作品}.
\end{equation}
同樣地，這個主張在講座裡是局部的，而不是形上學上的普遍真理；但它卻是整場訪談邏輯上的中心。時間的代價之所以合理，正是因為最後得到的書確實更好。

在這裡，再補上一張逐字稿式的示意圖會很有幫助：

\begin{figure}[h]
\centering
\begin{tikzpicture}
\node[draw, rectangle, minimum width=3.2cm, minimum height=1cm] (P) at (0,0) {規劃：1 年};
\node[draw, rectangle, minimum width=3.2cm, minimum height=1cm] (F) at (4,0) {第一稿：1 年};
\node[draw, rectangle, minimum width=3.2cm, minimum height=1cm] (S) at (8,0) {第二稿：1 年};
\draw[->] (P) -- (F);
\draw[->] (F) -- (S);
\end{tikzpicture}
\end{figure}

這張圖的目的並不是裝飾，而是明白顯示：在這場講座裡，品質被視為分散在各個階段中完成的，而不是集中於某一瞬間的靈光一現。

\section{中途襲來的嚴重懷疑，以及繼續下去的規則}

講座在這裡忽然轉入內在層面。訪談者問到焦慮、寫作阻滯，以及崩潰時刻。Follett 沒有用那種「受苦靈魂」的大詞，但也沒有否認困難。他只是點出一個每本書都會出現的時刻：人已經寫到很深的地方，累積了幾百頁，然後突然冒出一個念頭——到底會有誰想讀這個？

這個時刻可以表示為
\begin{equation}
\text{已寫出數百頁} \Rightarrow \text{嚴重懷疑}.
\end{equation}
而講座真正的力量，在於接下來會發生什麼。答案不是啟示，不是自我表達，也不是恐慌，而是操作性的：
\begin{align}
\text{懷疑} &\Rightarrow \text{繼續往前} \Rightarrow \text{若有需要，之後再重寫},\\
\text{如果稿子還不夠好} &\Rightarrow \text{重寫} \Rightarrow \text{改進}.
\end{align}
這大概是整場講座最實際的一條定理。一次心理崩塌，是藉由對程序的信任而被解決的。人不在那一刻停下來，任由那種感覺成為最高權威，而是繼續走下去，因為他知道修訂始終存在。

\subsection{Question \& Answer}

\paragraph{Question.} 當一本書忽然連作者自己都覺得根本看不下去時，應該怎麼辦？

\paragraph{Answer.} 講座的答案故意不帶戲劇性：繼續。先把這個念頭放到一邊，完成眼前的工作，並把最後的判決交給後面的修訂。那種「不可讀」的感覺，不被授予最終裁判權。

\section{標準、投入，以及對工作的熱愛}

講座本可以在懷疑與忍耐中結束，但它沒有。它選擇了更強的一個收束方式。Follett 回憶起曾有人告訴他，他作為作家的問題在於：他不是一個受苦的靈魂。他熱愛這個過程，他沉浸其中。這個最後的關係並不瑣碎：
\begin{equation}
\text{對工作的投入} \Rightarrow \text{對過程的忍受力}.
\end{equation}
這一點替前面的一切做了澄清。三個螢幕、整本重寫、三年的結構、以及在中途崩潰時仍然不停止——這一切之所以能持續，不是因為工作只是被勉強忍受，而是因為工作本身值得被熱愛。

這也讓講座回到它那個隱藏的開場對照。真正的魔法從來不來自房間，而來自方法。方法之所以還能被活下去，則是因為它被依戀所驅動，而不是靠陰鬱的自我懲罰所維持。

\section{Summary}

這場講座是一路把各種浪漫想像一一轉化成方法。它先從那間書房開始，隨即拒絕把房間當作祕密。接著，它轉入三螢幕的工作流，使起草、回饋與重寫被組織成一個可見的過程。一條來自讀者的評論，隨後被翻譯成節奏上的診斷，再拿去對照「每三到四頁要有一個轉折」的經驗法則。然後，修訂被重新定義：第二稿不是加了註記的第一稿，而是整本從頭重寫的新書。這條標準的算術形式，則在規劃一年、第一稿一年、第二稿一年這三段式分解中被明確寫了出來。最後，講座面對那種寫到幾百頁後必然出現的嚴重懷疑，並用一條「先繼續、之後再讓修訂來裁決」的規則來作答。講座最後的尾音，與它的方法本身同樣重要：嚴格之所以可以被持續，正因為工作是被熱愛的。